% Part: lambda-calculus
% Chapter: lambda-definability
% Section: lambda-definable-recursive

\documentclass[../../../include/open-logic-section]{subfiles}

\begin{document}

\olfileid[fr]{lam}{rep}{ldr}
\olsection{Les fonctions \texorpdfstring{$\lambda$}{λ}-définissables sont récursives}

Toutes les fonctions récursives
partielles sont $\lambda$-définissables,
et la réciproque est également vraie :
toute fonction $\lambda$-définissable
est récursive partielle.

\begin{thm}
\ollabel{thm:lambda-computable}
Si une fonction partielle~$f$ est
$\lambda$-définissable, alors elle
est récursive partielle.
\end{thm}

\begin{proof}
Esquissons seulement la preuve.
Commençons par arithmétiser les
lambda-termes : attribuons-leur
systématiquement des numéros de Gödel,
en utilisant le codage usuel des suites
par les puissances de nombres premiers.
Définissons ensuite une fonction
récursive partielle $\fn{normalize}(t)$
qui reçoit le numéro de Gödel~$t$
d'un lambda-terme et renvoie celui
de sa forme normale si elle existe;
sinon, elle est indéfinie. Une recherche
effective de réduction en ordre normal
fournit cette fonction, en utilisant le
fait qu'un terme ayant une forme normale
est normalisé par cette stratégie.
Définissons aussi deux fonctions
récursives partielles :
$\fn{toChurch}$ envoie un entier naturel
sur le numéro de Gödel du numéral de
Church correspondant, et
$\fn{fromChurch}$ effectue l'opération
inverse sur les codes de ces numéraux.

Soit maintenant~$F$ un lambda-terme
qui définit~$f$. Pour calculer
$f(n_1,\dots,n_k)$, obtenons d'abord
les codes des numéraux correspondants
au moyen de $\fn{toChurch}(n_i)$.
À partir de ces codes et de~$\Gn{F}$,
construisons effectivement le code
du terme $F\num{n_1}\dots\num{n_k}$.
Appliquons-lui ensuite
$\fn{normalize}$. Si $f(n_1,\dots,n_k)$
est définie, ce terme a pour forme
normale un numéral de Church; sinon,
il n'a pas de forme normale et
\[
  \fn{normalize}
    (\Gn{F\num{n_1}\dots\num{n_k}})
\]
est indéfinie. Dans le premier cas,
$\fn{fromChurch}$ appliquée au code
obtenu rend la valeur~$f(n_1,\dots,n_k)$.
Cette composition de fonctions
récursives partielles calcule donc~$f$.
\end{proof}

La procédure utilise la stratégie externe gauche pour calculer
la forme normale lorsqu'elle existe. La standardisation est
énoncée dans la preuve du
\olref[lam][rep][par]{thm:partial-recursive-lambda}.
Pour reconnaître un numéral et en lire l'entier, on peut employer
les indices de de Bruijn, qui rendent la reconnaissance
indépendante des noms des variables liées.

\end{document}
